\chapter{CTL over Generated Futures}
\label{ch:ctl}

\index{CTL}
\index{temporal logic}
A grammar generates a branching transition system. CTL speaks directly about that
branching: path quantifiers distinguish \emph{there exists a future} from
\emph{every future}, while temporal operators distinguish the next state,
eventuality, invariance, and until.

\begin{center}
\begin{tabular}{@{}lll@{}}
\toprule
Formula & Reading & Typical use \\
\midrule
$EF\,p$ & some future can reach $p$ & possibility / witness \\
$AF\,p$ & every future eventually reaches $p$ & termination / response \\
$EG\,p$ & some path preserves $p$ forever & viable invariant path \\
$AG\,p$ & every reachable state satisfies $p$ & safety invariant \\
$(p\ AU\ q)$ & every path keeps $p$ until $q$ & constrained progress \\
$(p\ EU\ q)$ & some path keeps $p$ until $q$ & possible progress \\
\bottomrule
\end{tabular}
\end{center}

The path quantifier is part of the until operator. We write \texttt{AU} and
\texttt{EU} as binary algebraic operators, not an outer \texttt{A} or \texttt{E}
wrapped around a bracketed until formula. Consequently a subexpression never
changes meaning merely because another expression surrounds it:

\begin{lstlisting}
def boundedUntilTerminal : CTL Observation :=
  atom withinCapacity AU atom isTerminal

def someAuthorizedContinuation : CTL Observation :=
  atom hasValidAuth EU atom isSucceeded
\end{lstlisting}

This convention also makes tree structure explicit. In
\texttt{r \&\& (p AU q)}, the right operand is the complete \texttt{AU} expression;
there is no context-sensitive reinterpretation of \texttt{[p U q]}.

For a request event carrying identifier $r$, an expected response property might
say that on every continuation a matching success or failure eventually occurs.
Authentication might say that every successful use of a session has an earlier
issuance event with the same session identifier.

\section{Past claims}

\index{past-time logic}
CTL is future-facing at the current state, but security and protocol claims often
refer to evidence in the past: ``this capability was previously issued'' or ``this
response matches an earlier request.'' Since states are trace prefixes, past facts
can initially be expressed as predicates over the accumulated trace. This gives a
clear baseline without prematurely committing the surface language to a particular
past-time temporal logic.

There is still an open design question. Explicit past operators may make properties
more readable and allow better checking algorithms. The project should compare
trace predicates, CTL with past operators, and compilation of finite-history
monitors into state. The criterion is maintainability of specifications as the
event schema evolves.

\section{Fairness and time bounds}

$AF$ does not by itself distinguish a participant that is perpetually denied a
turn from one that violates the protocol. Fairness assumptions must therefore be
named, scoped, and visible in theorem statements. Similarly, ``eventually'' is not
``within 200 milliseconds.'' Quantitative bounds belong in an explicit timed layer
rather than being smuggled into timestamps and informal prose.
